\chapter{Desenvolvimento e Implementação}
\label{cap:implementacao}

Este capítulo documenta a engenharia de software da solução desenvolvida, descrevendo a integração entre o motor gráfico Unity, o middleware Neo Cortex e o ambiente de inferência local Ollama. Detalha-se a implementação da interface de usuário, o protocolo de comunicação cliente-servidor, os algoritmos de triagem baseados em ReAct e as especificações da infraestrutura computacional utilizada.

\section{Visão Geral da Arquitetura do Sistema}
\label{sec:dev_arquitetura}

Para viabilizar a simulação de personagens autônomos sem dependência de serviços externos em nuvem, concebeu-se uma arquitetura em camadas desacopladas. A Figura \ref{fig:arquitetura_pipeline} sintetiza o fluxo de dados entre as camadas do sistema:

\begin{figure}[htbp]
\centering
\begin{tikzpicture}[
    auto,
    node distance=1.0cm and 1.2cm,
    box/.style={rectangle, draw=black!75, rounded corners=3pt, thick, align=center, font=\footnotesize, minimum height=0.85cm},
    client/.style={box, fill=blue!10, draw=blue!80, minimum width=3.2cm},
    mid/.style={box, fill=purple!10, draw=purple!80, minimum width=3.2cm},
    db/.style={cylinder, draw=orange!80, fill=orange!10, shape border rotate=90, aspect=0.25, align=center, font=\scriptsize, minimum width=2.4cm, minimum height=0.9cm},
    llm/.style={box, fill=green!15, draw=green!60!black, minimum width=3.6cm, font=\footnotesize\bfseries},
    arr/.style={thick, ->, >=stealth}
]

    % Camada 1: Cliente e Apresentação
    \node (player) [client, fill=yellow!15, draw=orange!80] {\textbf{Jogador (Detetive)}\\Inquérito em linguagem natural};
    \node (unity) [client, right=1.2cm of player] {\textbf{Unity Engine (C\#)}\\Interface de Diálogo e Máquina de Estados};

    % Camada 2: Mediação e Triagem
    \node (neocortex) [mid, below=1.1cm of unity] {\textbf{Neo Cortex Middleware}\\Tratamento e Despachador HTTP};
    \node (triage) [mid, fill=purple!20, below=1.1cm of player] {\textbf{Agente de Triagem (Gemma-2)}\\Ciclo ReAct: Decisão de Contexto};

    % Camada 3: Base de Dados e Chunking
    \node (dbnode) [db, below=1.0cm of triage] {\textbf{Base Factual JSON}\\Perfil, Cronograma, Relações};
    \node (chunker) [box, fill=orange!15, draw=orange!80, right=1.2cm of dbnode, minimum width=3.2cm] {\textbf{Filtro de Chunking}\\Fragmentação Temporal Dinâmica};

    % Camada 4: Inferência do Modelo Principal
    \node (generator) [llm, below=1.0cm of chunker] {\textbf{Modelo da Persona (LLM Local)}\\Gemma-2 / Llama-3.1 / Qwen-2.5 / DeepSeek};

    % Conexões do Fluxo
    \draw [arr] (player) -- node [above, font=\scriptsize] {1. Pergunta livre} (unity);
    \draw [arr] (unity) -- node [right, font=\scriptsize] {2. Requisição} (neocortex);
    \draw [arr] (neocortex) -- node [above, font=\scriptsize] {3. Inquérito bruto} (triage);
    \draw [arr] (triage) -- node [left, font=\scriptsize] {4. Comando [BUSCAR\_*]} (dbnode);
    \draw [arr] (dbnode) -- node [above, font=\scriptsize] {5. Extração} (chunker);
    \draw [arr] (chunker) -- node [right, font=\scriptsize] {6. Contexto Filtrado Mínimo} (generator);
    
    % Retorno ao Unity
    \draw [arr] (generator.east) -- ++(0.7,0) |- node [pos=0.25, right, font=\scriptsize] {7. Réplica do NPC} (unity.east);

    % Caixa de delimitação do ambiente local (Ollama)
    \begin{scope}[on background layer]
        \node (localbox) [draw=gray!50, dashed, fill=gray!4, rounded corners=6pt, fit=(triage) (dbnode) (chunker) (generator), label={[yshift=-0.5cm]above:\scriptsize\textbf{Ambiente de Execução Local (Ollama Runtime / Hardware do Usuário)}}] {};
    \end{scope}

\end{tikzpicture}
\caption{Arquitetura em camadas do pipeline Agentic RAG integrado ao motor Unity via Neo Cortex.}
\label{fig:arquitetura_pipeline}
\end{figure}


O sistema organiza-se em três domínios principais:
\begin{enumerate}
    \item \textbf{Camada de Apresentação e Interação (Unity Engine):} Gerencia a renderização gráfica, o estado do jogo e a captura das mensagens digitadas pelo usuário;
    \item \textbf{Camada de Mediação e Orquestração (Neo Cortex Middleware):} Atua como a ponte assíncrona entre o Unity e os modelos de linguagem, executando o ciclo de triagem ReAct e o despachador de fragmentos JSON;
    \item \textbf{Camada de Inferência Neural Local (Ollama Runtime):} Responsável pela alocação dos pesos dos modelos quantizados em memória VRAM e pela execução do cálculo tensorial de geração de \textit{tokens}.
\end{enumerate}

\section{Implementação no Motor Gráfico Unity}
\label{sec:dev_unity}

O ambiente visual do jogo de investigação foi desenvolvido utilizando o motor gráfico Unity \textcolor{red}{[TODO: faltando informação: especificar a versão exata do Unity utilizada, e.g., Unity 2022.3.x LTS]}. 

A cena principal simula a sala de interrogatório de uma delegacia ou a biblioteca da mansão onde os suspeitos aguardam o inquérito. A interface de usuário (\textit{User Interface}, UI) foi estruturada através dos seguintes elementos gráficos:

\begin{itemize}
    \item \textbf{Área de Histórico de Diálogo (\textit{Scroll View} com TextMeshPro):} Exibe cronologicamente as intervenções do detetive e as respostas geradas pelo NPC, com suporte a cores diferenciadas e rubricas teatrais em itálico;
    \item \textbf{Campo de Entrada de Texto (\textit{TMP\_InputField}):} Permite a digitação irrestrita de perguntas pelo jogador, enviadas ao pressionar a tecla \textit{Enter} ou o botão de confirmação;
    \item \textbf{Indicador de Estado / \textit{Loading Spinner}:} Elemento visual que sinaliza quando o agente de triagem ou o modelo gerador está processando o turno, fornecendo \textit{feedback} imediato ao jogador para mascarar a latência de inferência local.
\end{itemize}

O script principal em C\#, denominado \texttt{InvestigationDialogueManager.cs}, coordena a máquina de estados da interface. A listagem de código a seguir ilustra a rotina assíncrona utilizada para despachar a entrada do jogador sem travar a taxa de quadros (\textit{framerate}) da engine:

\begin{lstlisting}[language={[Sharp]C}, caption={Rotina assíncrona em C\# para comunicação com o middleware Neo Cortex.}, label={lst:envio_mensagem}]
public IEnumerator SendPlayerInquiry(string playerText) {
    SetUIState(UIState.WaitingForResponse);
    AppendMessageToHistory("Detetive", playerText);

    // Estrutura payload JSON para o middleware Neo Cortex
    InquiryPayload payload = new InquiryPayload {
        characterId = "alberto_castelo",
        userInput = playerText,
        currentInGameTime = "23:30"
    };

    string jsonPayload = JsonUtility.ToJson(payload);
    using (UnityWebRequest request = new UnityWebRequest(
               "http://localhost:11434/api/generate", "POST")) {
        byte[] bodyRaw = Encoding.UTF8.GetBytes(jsonPayload);
        request.uploadHandler = new UploadHandlerRaw(bodyRaw);
        request.downloadHandler = new DownloadHandlerBuffer();
        request.SetRequestHeader("Content-Type", "application/json");

        yield return request.SendWebRequest();

        if (request.result == UnityWebRequest.Result.Success) {
            ProcessServerResponse(request.downloadHandler.text);
        } else {
            HandleNetworkError(request.error);
        }
    }
    SetUIState(UIState.ReadyForInput);
}
\end{lstlisting}

\section{O Middleware Neo Cortex e o Protocolo de Mensagens}
\label{sec:dev_neocortex}

O framework Neo Cortex atua como um barramento de integração que abstrai as chamadas brutas de API do Ollama e padroniza a conversão de tipos de dados entre o C\# do Unity e as cadeias de texto dos LLMs.

O protocolo de comunicação baseia-se em requisições assíncronas HTTP REST direcionadas em malha fechada ao endereço local (\texttt{127.0.0.1}). A escolha por chamadas locais garante mínima sobrecarga de rede (latência de transporte inferior a 2 milissegundos), assegurando que o tempo de resposta percebido pelo jogador seja dominado exclusivamente pelo tempo de inferência dos modelos.

\section{O Algoritmo de Triagem ReAct e os Comandos Estruturados}
\label{sec:dev_triagem}

Conforme delineado no modelo conceitual da Seção \ref{sec:met_arquitetura_rag}, a viabilidade do \textit{Agentic RAG} depende da capacidade do agente de triagem em interpretar a intenção interrogativa do jogador e despachar exclusivamente a partição documental correspondente. Nesta seção, detalha-se a implementação desse mecanismo na camada de engenharia, especificando a parametrização do modelo, o ciclo de raciocínio ReAct e a gramática formal de comandos estruturados.

Utilizando o modelo Gemma-2-9b-it configurado com diretrizes rigorosas de sistema, o agente recebe a pergunta bruta do jogador e avalia se há necessidade de consulta à base factual antes de invocar o modelo do personagem. O \textit{prompt} de sistema do agente de triagem impõe que a resposta seja emitida estritamente sob o formato padronizado de pensamento e ação:

\begin{quote}
\small
\texttt{Thought: [raciocínio analítico identificando intenção e período de tempo]}\\
\texttt{Action: [comando estruturado selecionado]}
\end{quote}

A Tabela \ref{tab:comandos_triagem} detalha o mapa de correspondência entre as intenções do interrogador, as ações emitidas pelo agente e os arquivos JSON despachados em memória:

\begin{table}[htbp]
\centering
\caption{Gramática de Ações do Agente de Triagem e Despacho de Contexto}
\vspace{0.2cm}
\label{tab:comandos_triagem}
\footnotesize
\renewcommand{\arraystretch}{1.3}
\begin{tabularx}{\textwidth}{|>{\hsize=0.95\hsize}X|>{\hsize=1.15\hsize}X|>{\hsize=0.9\hsize}X|}
\hline
\textbf{Intenção da Pergunta} & \textbf{Comando Emitido (\textit{Action})} & \textbf{Fragmento JSON Injetado} \\ \hline
Perguntas sobre ações entre 08:00 e 12:00 & {\scriptsize\texttt{[BUSCAR\_CRONOGRAMA\_MANHA]}} & {\scriptsize\texttt{cronograma\_manha.json}} \\ \hline
Perguntas sobre ações entre 12:00 e 18:00 & {\scriptsize\texttt{[BUSCAR\_CRONOGRAMA\_TARDE]}} & {\scriptsize\texttt{cronograma\_tarde.json}} \\ \hline
Perguntas sobre ações após as 18:00 até a madrugada & {\scriptsize\texttt{[BUSCAR\_CRONOGRAMA\_NOITE]}} & {\scriptsize\texttt{cronograma\_noite.json}} \\ \hline
Perguntas sobre opiniões ou suspeitas contra terceiros & {\scriptsize\texttt{[BUSCAR\_RELACIONAMENTOS]}} & {\scriptsize\texttt{relacionamentos.json}} \\ \hline
Saudações, perguntas sobre biografia ou provocações genéricas & {\scriptsize\texttt{[RESPONDER\_DIRETO]}} & Nenhum (apenas \texttt{perfil.json} base) \\ \hline
\end{tabularx}
\end{table}

Após a identificação do comando pela expressão regular do despachador, o conteúdo do fragmento JSON é concatenado ao \textit{prompt} de persona do modelo principal, garantindo que o volume de \textit{tokens} transmitido permaneça na faixa ótima de processamento.

\section{Infraestrutura Computacional e Ambiente de Execução Local}
\label{sec:dev_hardware}

A validação experimental foi conduzida integralmente em ambiente local. A Tabela \ref{tab:setup_hardware} sintetiza as especificações do ambiente de teste \textcolor{red}{[TODO: validar se as especificações de hardware abaixo coincidem exatamente com a máquina onde os testes foram executados]}.

\begin{table}[htbp]
\centering
\caption{Especificações do Ambiente de Execução Local}
\vspace{0.2cm}
\label{tab:setup_hardware}
\small
\renewcommand{\arraystretch}{1.3}
\begin{tabularx}{\textwidth}{|p{4.5cm}|X|}
\hline
\textbf{Componente} & \textbf{Especificação Técnica} \\ \hline
Sistema Operacional & Linux (Ubuntu 24.04 LTS / Debian-based) / Windows 11 \\ \hline
Processador (CPU) & \textcolor{red}{[faltando informação: e.g., AMD Ryzen 7 / Intel Core i7]} \\ \hline
Placa Gráfica (GPU) & \textcolor{red}{[faltando informação: e.g., NVIDIA GeForce RTX 3060 / 4060, 12GB VRAM]} \\ \hline
Memória RAM do Sistema & 16 GB / 32 GB DDR4/DDR5 \\ \hline
Orquestrador de LLMs & Ollama (v0.3.x+) \\ \hline
Formato de Quantização & GGUF Q4\_K\_M (4 bits) \\ \hline
Temperatura Fixa & 0.1 \\ \hline
\end{tabularx}
\end{table}

A utilização da quantização Q4\_K\_M permitiu que os modelos de 7B a 9B parâmetros ocupassem entre 4,8 GB e 6,2 GB de VRAM, permitindo a execução simultânea do motor gráfico Unity e do servidor Ollama na mesma máquina sem estouro de memória gráfica (\textit{Out of Memory}, OOM).
